\documentclass[11pt,a4paper]{article}

\usepackage[margin=2.2cm]{geometry}
\usepackage[T1]{fontenc}
\usepackage[utf8]{inputenc}
\usepackage{lmodern}
\usepackage{booktabs}
\usepackage{array}
\usepackage{enumitem}
\usepackage{amsmath}
\usepackage{xcolor}
\usepackage{fancyhdr}
\usepackage{lastpage}
\usepackage{hyperref}
\usepackage{multicol}

\hypersetup{
  colorlinks=true,
  linkcolor=black,
  urlcolor=blue
}

\pagestyle{fancy}
\fancyhf{}
\lhead{COMP90042 NLP}
\rhead{Practice Exam}
\cfoot{Page \thepage\ of \pageref{LastPage}}

\setlength{\parindent}{0pt}
\setlength{\parskip}{0.55em}
\setlist[itemize]{topsep=0.25em,itemsep=0.15em}
\setlist[enumerate]{topsep=0.35em,itemsep=0.45em}

\newcommand{\qmarks}[1]{\hfill\textbf{[#1 marks]}}
\newcommand{\questionmeta}[3]{\textit{Topic: #1. Estimated time: #2. Difficulty: #3.}}

\begin{document}

\begin{center}
{\Large \textbf{COMP90042 Natural Language Processing}}\\
{\large Practice Exam Paper}\\[0.4em]
Generated from \texttt{l23-review-part1-v1.pdf} using the \texttt{practice-exams} skill
\end{center}

\section*{Exam Information}

\begin{tabular}{@{}ll@{}}
\toprule
Format & On-campus, closed book \\
Allowed aid & Non-programmable calculator \\
Writing time & 120 minutes, plus 15 minutes reading time \\
Total marks & 120 \\
Structure & Part A short answer, Part B method questions, Part C algorithm questions \\
\bottomrule
\end{tabular}

\section*{Coverage Blueprint}

\begin{tabular}{@{}p{0.35\linewidth}r p{0.43\linewidth}@{}}
\toprule
\textbf{Area} & \textbf{Marks} & \textbf{Main skills tested} \\
\midrule
Preprocessing and tokenisation & 10 & Definitions, design choices, examples \\
N-gram language models & 20 & Derivation, smoothing, interpolation \\
Text classification & 15 & Model comparison, bias-variance, feature design \\
POS tagging and HMMs & 25 & Tagging models, emissions/transitions, Viterbi \\
Neural networks for NLP & 20 & FFN/RNN/LSTM design and interpretation \\
Transformers and pretrained models & 20 & Attention, model architecture, GPT/BERT/T5 \\
Word embeddings & 10 & Distributional matrices, PMI/PPMI, evaluation \\
\bottomrule
\end{tabular}

\section*{Part A: Short Answer Questions \hfill 36 marks}

Answer each question in 2--3 sentences. Each question is worth 3 marks.

\begin{enumerate}
\item Explain the difference between sentence segmentation and tokenisation. Give one example where tokenisation is not straightforward.
\item What is the difference between derivational and inflectional morphology, and why might the distinction matter for word normalisation?
\item Contrast lemmatisation and stemming. Give one advantage and one disadvantage of each.
\item In an n-gram language model, what independence assumption is being made?
\item Why is smoothing necessary in n-gram language modelling?
\item What is the purpose of interpolation in language modelling?
\item Give two examples of text classification tasks and explain what the input and output are in each.
\item Explain the bias-variance trade-off in the context of text classification.
\item What is the difference between open-class and closed-class POS categories?
\item In an HMM POS tagger, what do emission probabilities and transition probabilities represent?
\item Why are word embeddings useful compared with one-hot word representations?
\item What is the main architectural difference between GPT and BERT?
\end{enumerate}

\section*{Part B: Method Questions \hfill 60 marks}

Answer all questions. Show reasoning clearly.

\subsection*{B1. Preprocessing Pipeline Design \qmarks{12}}
\questionmeta{Preprocessing and classification}{12 minutes}{medium}

You are building a sentiment classifier for short social media posts. Describe a preprocessing pipeline that includes sentence segmentation, tokenisation, subword handling, word normalisation, and stop word handling. For each step, explain one design choice and one possible risk.

\subsection*{B2. Smoothing Methods \qmarks{16}}
\questionmeta{N-gram language models}{16 minutes}{hard}

Compare add-\(k\) smoothing, absolute discounting, Katz backoff, Kneser-Ney smoothing, and interpolation. Your answer should explain the core motivation of each method and identify one situation where a more sophisticated method is preferable to add-\(k\) smoothing.

\subsection*{B3. Classifier Choice \qmarks{16}}
\questionmeta{Text classification}{16 minutes}{hard}

A course wants to classify student discussion posts by topic. Compare Naive Bayes, logistic regression, SVM, kNN, and a neural network for this task. Discuss assumptions, data requirements, interpretability, and likely bias-variance behaviour.

\subsection*{B4. Pretrained Model Selection \qmarks{16}}
\questionmeta{Transformers and pretrained models}{16 minutes}{medium}

For each task below, choose GPT, BERT, or T5 and justify your choice: sentiment classification, open-ended story generation, and rewriting a sentence into a simpler form. Include one limitation or risk for each selected model.

\section*{Part C: Algorithm Questions \hfill 24 marks}

Answer both questions. Show calculations or algorithm steps.

\subsection*{C1. N-gram Probability and Interpolation \qmarks{12}}
\questionmeta{Language modelling}{12 minutes}{medium}

Suppose a corpus gives the following counts:
\begin{itemize}
\item \(\mathrm{count}(\text{the cat}) = 8\), \(\mathrm{count}(\text{the}) = 20\)
\item \(\mathrm{count}(\text{cat sat}) = 3\), \(\mathrm{count}(\text{cat}) = 10\)
\item \(\mathrm{count}(\text{sat down}) = 1\), \(\mathrm{count}(\text{sat}) = 5\)
\item unigram probabilities: \(P(\text{cat})=0.04\), \(P(\text{sat})=0.02\), \(P(\text{down})=0.01\)
\end{itemize}

Estimate the interpolated probability of \textit{the cat sat down} using bigram and unigram models with \(\lambda_{\text{bigram}} = 0.7\) and \(\lambda_{\text{unigram}} = 0.3\). Use:
\[
P(w_i \mid w_{i-1}) = 0.7P_{\text{bigram}}(w_i \mid w_{i-1}) + 0.3P_{\text{unigram}}(w_i)
\]

\subsection*{C2. One Viterbi Step \qmarks{12}}
\questionmeta{HMM POS tagging}{12 minutes}{medium}

Consider a tiny HMM POS tagger with tags \(N\) and \(V\). For the sentence \textit{fish swim}, use the following values:
\begin{itemize}
\item Initial probabilities: \(P(N)=0.6\), \(P(V)=0.4\)
\item Emissions: \(P(\text{fish}\mid N)=0.5\), \(P(\text{fish}\mid V)=0.2\), \(P(\text{swim}\mid N)=0.1\), \(P(\text{swim}\mid V)=0.6\)
\item Transitions: \(P(N\mid N)=0.3\), \(P(V\mid N)=0.7\), \(P(N\mid V)=0.4\), \(P(V\mid V)=0.6\)
\end{itemize}

Run Viterbi for the two words and give the most likely tag sequence with its probability.

\clearpage
\section*{Final Answer Key}
\thispagestyle{fancy}
\enlargethispage{2\baselineskip}
\scriptsize
\setlength{\parskip}{0.18em}
\begin{multicols}{2}

\textbf{Part A}
\begin{enumerate}[leftmargin=1.2em,itemsep=0.15em]
\item Segmentation splits sentences; tokenisation splits tokens. Hard cases: abbreviations, contractions, URLs, hashtags, emojis, or no-whitespace languages.
\item Inflection changes grammatical form; derivation creates a new word/category. Normalisation should not always collapse derivational forms.
\item Lemmatisation uses linguistic analysis and dictionary forms; stemming heuristically strips affixes. Lemmas are cleaner; stems are faster but noisier.
\item The next word depends only on the previous \(n-1\) words.
\item Smoothing avoids zero probabilities and reserves probability mass for unseen n-grams.
\item Interpolation combines higher- and lower-order n-gram estimates for robustness.
\item Topic classification, sentiment analysis, or native language identification: input text, output class label.
\item High bias underfits; high variance overfits sparse text features, especially with little data.
\item Open classes admit new words, e.g. nouns/verbs; closed classes are relatively fixed, e.g. determiners/prepositions.
\item Emissions: \(P(\text{word}\mid\text{tag})\). Transitions: \(P(\text{tag}_i\mid\text{tag}_{i-1})\).
\item Embeddings are dense, similarity-aware representations; one-hot vectors are sparse and do not encode semantic closeness.
\item GPT is decoder-only/autoregressive for generation; BERT is encoder-only/masked-LM for contextual representation.
\end{enumerate}

\textbf{Part B}

\textbf{B1.} Suitable pipeline: sentence segmentation, social-media-aware tokenisation, subword handling for rare/OOV words, cautious lemmatisation/stemming, and selective stop-word removal. Main risks: losing negation, emojis, hashtags, domain terms, or sentiment cues.

\textbf{B2.} Add-\(k\): simple count adjustment but often over-smooths. Absolute discounting: subtracts fixed mass from seen n-grams. Katz: discounted seen n-grams plus backoff for unseen ones. Kneser-Ney: uses continuation probability, good for sparse data. Interpolation: mixes orders directly.

\textbf{B3.} NB: simple/high-bias/independence assumption. Logistic regression: discriminative and interpretable. SVM: strong for sparse high-dimensional features. kNN: simple but distance-sensitive and costly. Neural networks: richer representations, more data/regularisation, higher variance.

\textbf{B4.} Sentiment: BERT, because encoder representations support classification. Story generation: GPT, because decoder-only autoregression supports continuation. Simplification: T5, because encoder-decoder text-to-text modelling supports rewriting. Mention risks such as hallucination, data mismatch, or limited controllability.

\textbf{Part C}

\textbf{C1.}
\[
P(c\mid t)=0.7(8/20)+0.3(0.04)=0.292
\]
\[
P(s\mid c)=0.7(3/10)+0.3(0.02)=0.216
\]
\[
P(d\mid s)=0.7(1/5)+0.3(0.01)=0.143
\]
Final probability, ignoring the first-word prior:
\[
0.292\times0.216\times0.143=0.009018816.
\]

\textbf{C2.} Initialise: \(N=0.6\times0.5=0.30\), \(V=0.4\times0.2=0.08\). For \textit{swim}: ending \(N=\max(0.009,0.0032)=0.009\); ending \(V=\max(0.126,0.0288)=0.126\). Best sequence: \(N,V\), probability \(0.126\).

\end{multicols}

\end{document}
